\begin{exsol}{22}{MWG Exercise 8.B.7}{262}{Advanced}{B2}
\asked Show that a strictly dominant strategy of the mixed extension of a game,
$[I,\{\Delta(S_i)\}, \{u_i(\cdot)\}]$, must be a pure strategy.

\method Mixed strategies and dominance in the mixed extension (MWG Def.~8.B.4, p.~240;
Prop.~8.B.1 and Exercise~8.B.6, p.~241). Key fact: expected utility is linear in one's own
mixing probabilities, so a mixture's payoff is a weighted average of the payoffs of the pure
strategies it uses. Mixed strategies are the subject of Lecture~2; only this averaging property
is needed here.

\solution
Fix player~$i$ and a mixed strategy $\sigma_i \in \Delta(S_i)$. $\sigma_i$ is \emph{strictly
dominant} if $u_i(\sigma_i, \sigma_{-i}) > u_i(\sigma_i', \sigma_{-i})$ for every
$\sigma_i' \ne \sigma_i$ and every $\sigma_{-i}$. In particular this must hold against every pure
profile $\sm$ (a degenerate $\sigma_{-i}$) and every pure $\sigma_i' = s_i \ne \sigma_i$.

\textbf{Proof.} Suppose $\sigma_i$ is strictly dominant but not pure: its support contains at
least two pure strategies $s_i^1, \dots, s_i^N$ ($N \ge 2$), with $\sigma_i(s_i^n) > 0$ and
$\sum_n \sigma_i(s_i^n) = 1$. Each $s_i^n$ is then a strategy different from $\sigma_i$. Fix any
$\sm \in \Sm$. By linearity of expected utility,
\[
  u_i(\sigma_i, \sm) = \sum_{n=1}^{N} \sigma_i(s_i^n)\, u_i(s_i^n, \sm)
  \le \max_{n} u_i(s_i^n, \sm) .
\]
Let $n^*$ attain the maximum. Then $u_i(\sigma_i, \sm) \le u_i(s_i^{n^*}, \sm)$, while strict
dominance of $\sigma_i$ requires $u_i(\sigma_i, \sm) > u_i(s_i^{n^*}, \sm)$. Contradiction. Hence
the support of $\sigma_i$ is a single pure strategy: $\sigma_i$ is pure. \hfill$\square$

\textbf{Remarks.}
\begin{itemize}
  \item The argument needs only \emph{one} opponent profile $\sm$; it uses no finiteness of
  $\Sm$ and works for any $I$.
  \item It fails for \emph{weak} dominance in the ``$\ge$'' sense. If two pure strategies tie
  against everything (as $U$ and $D$ in item~12), every mixture of them is also ``$\ge$-dominant''.
  \item The converse direction is different: a pure strategy can be strictly \emph{dominated} by a
  mixed strategy without being dominated by any pure one (MWG Fig.~8.B.5, pp.~240--241). Mixing can help
  in eliminating strategies, never in creating a dominant one.
\end{itemize}

\interp A strictly dominant strategy must be the unique best action against every opponent
profile. A mixture is never strictly better than the best of the pure strategies it mixes, so
mixing can only be optimal when the player is indifferent among those strategies, and
indifference is incompatible with strict dominance. This is the same reason why, in Lecture~2,
a player mixes only among strategies that give her equal payoffs.

\verify \texttt{check\_C.py} draws 20\,000 random supports (2--4 pure strategies, random payoffs
and weights, exact fractions) and confirms that the mixture's payoff never exceeds the best
support payoff, so it never strictly beats every pure strategy in its support. \emph{Comparison
with the MWG solutions manual} (p.~8-4): same argument (a strict inequality against each pure
strategy in the support, averaged with the mixing weights, contradicts linearity). It agrees.

\answer{If $\sigma_i$ mixes over $N \ge 2$ pure strategies, then for any $\sm$ its payoff is a
weighted average of theirs, so it is at most $\max_n u_i(s_i^n, \sm)$. It cannot strictly beat
that best pure strategy, which contradicts strict dominance. Hence every strictly dominant
strategy is pure.}
\end{exsol}
